\subsection{Why Sparse-Concept Calibration Vulnerability Is Dataset-Dependent}
\label{sec:a5_dataset_dependent}

Section~\ref{sec:a4_confounding} showed that $\log(1+f_{\mathrm{train}})$ remains independently associated with ECE after covariate adjustment, but that the strength of this association---and the visibility of a stratum-level ECE gradient---varies across datasets.
The remaining question is not whether ``datasets differ,'' but whether those differences correspond to measurable structural conditions.
Table~\ref{tab:a5_dataset_conditions} reports six such conditions, all computed from training-fold frequencies, training-only KC covariates, and the event-level SimpleKT ECE already reported in Table~\ref{tab:calibration_learner}.
Each cell is a FINDING (a structural or metric observation).
Explanations that would require KC ontologies, curriculum graphs, or item-text features---hierarchical placement, ceiling effects, item-feature richness, and tagging granularity---are labelled HYPOTHESIS and are not treated as established.

\subsubsection*{ASSISTments 2012: why a calibration gradient is observable}
Under the learner-based split, 18.9\% of ASSISTments KCs fall in sparse, very-sparse, or strict buckets, and the sparse stratum accumulates $N=403$ test events (Limited flag).
Event-level SimpleKT ECE increases monotonically from $0.113$ (dense) to $0.158$ (medium) to $0.225$ (sparse).
Two training-only couplings have the direction that would be expected if rarer concepts were also harder and later: frequency--difficulty $\rho=-0.227$ ($p<0.001$) and frequency--curriculum-position $\rho=-0.308$ ($p<0.001$).
The frequency--ECE association is the strongest of the three datasets ($\bar{\rho}=-0.60$).
These are jointly the conditions under which a sparse-KC calibration gradient is \emph{estimable and visible} in this dataset.

They do not, however, identify a single cause.
Item-support imbalance (median 205 items/KC in dense versus 1 in sparse) is partly mechanical: a KC with $f_{\mathrm{train}}<100$ cannot host hundreds of training items.
The non-mechanical item fact is that ASSISTments KCs are item-rich overall (median $44.5$ items/KC; frequency--item $\rho=0.92$).
\textbf{HYPOTHESIS:} the relatively small KC inventory ($n=265$) and right-skewed frequencies may reflect a coarse expert-tagged skill ontology in which low-frequency labels are genuinely rare or advanced skills.
Confirming that claim would require external KC metadata, which we do not use.

\subsubsection*{Junyi Academy: why learner-based sparse strata are inactive}
Junyi has $1{,}326$ KCs, but $97.5\%$ are dense and $0.0\%$ are sparse or very sparse under the pre-registered thresholds $\{20,100,500\}$.
Median training frequency is $7{,}074$---about seven times the ASSISTments median ($988$).
Consequently the learner-based sparse and very-sparse ECE cells are empty, not because calibration is known to be stable, but because the protocol's buckets receive no mass.
This is a FINDING about volume and thresholds, not about the absence of hard-to-calibrate concepts.
The same protocol populates sparse-like groups under the temporal split ($10$ KCs with $f_{\mathrm{train}}<100$, and a SimpleKT sparse ECE of $0.1624$ versus dense $0.0889$ in Table~\ref{tab:calibration_temporal}).

Where Junyi \emph{does} have contrast (dense versus medium), the direction matches ASSISTments: medium ECE is higher ($\Delta\mathrm{ECE}=+0.031$), frequency--difficulty coupling is present and in the expected direction ($\rho=-0.416$), and the frequency--ECE association is moderate ($\bar{\rho}=-0.43$).
Junyi item support is unusually uniform (median $18$ items/KC, IQR $5$; frequency--item $\rho=0.026$, $p=0.34$), so item-count heterogeneity is not a plausible driver of the medium-stratum ECE rise.
\textbf{HYPOTHESIS:} raising the dense threshold toward the Junyi frequency scale (for example $f_{\mathrm{train}}\ge 2000$) might reopen a sparse-like contrast under learner-based splits.
That is a threshold-sensitivity question, not a result of this study.

\subsubsection*{XES3G5M: why the pattern is not a second ASSISTments}
XES3G5M rules out the two simplest explanations for a missing gradient.
It has \emph{more} sparse-like mass than ASSISTments ($22.5\%$ of KCs) and \emph{better} sparse test support (sparse $N=2{,}002$, Reliable; very-sparse $N=109$, Limited).
A flat or inverted pattern here is therefore not an artefact of empty buckets or of $N<100$ instability.

What \emph{is} measurable and different:
(i)~SimpleKT event-level ECE is essentially flat from dense to sparse ($0.115\to 0.112\to 0.124$), while IRT is inverted ($0.154\to 0.107$ dense-to-sparse) and DKT increases---so ``the XES3G5M gradient'' is model-heterogeneous, not a single dataset law;
(ii)~the frequency--ECE association is the weakest of the three datasets ($\bar{\rho}=-0.23$);
(iii)~frequency--difficulty coupling is weak and opposite-signed ($\rho=+0.087$, $p=0.012$), not a strong ``sparse KCs are easier'' effect;
(iv)~KCs are easier overall (difficulty-proxy median $0.192$ versus $0.338$ on both other datasets);
(v)~item support is thin (median $3$ items/KC) and interaction mass is concentrated (top $10\%$ of KCs hold $67.9\%$ of training interactions; the single most frequent KC holds $19.4\%$).

\textbf{HYPOTHESIS:} a hierarchical KC tree, multi-skill expansion, a ceiling on advanced nodes, or item-side features could produce this combination.
None of those mechanisms is identified by the covariates we can construct from the processed interaction schema, so we do not promote them to findings.

\subsubsection*{Conditional conclusion}
Across the three datasets, an \emph{observable} SimpleKT sparse-KC calibration gradient coincides with three measurable conditions: (1)~non-trivial sparse mass under the split and thresholds actually used; (2)~enough sparse test events to support at least a Limited-flag ECE; and (3)~a frequency--difficulty coupling that is not inverted.
ASSISTments 2012 meets all three and shows a monotonic ECE rise.
Junyi Academy fails (1) and (2) on the learner-based split, so the protocol correctly withholds a sparse ECE rather than imputing stability.
XES3G5M meets (1) and (2) but not (3), and the SimpleKT gradient is flat; model-specific counter-patterns (inverted IRT, increasing DKT) further forbid a universal sparsity-to-miscalibration law.
These are diagnostic pre-conditions that future studies can check with the same protocol.
They are not a claim that training frequency causes miscalibration, nor that hierarchical curricula, ceiling effects, or tagging granularity have been demonstrated.
